% \section{Appendix B: Datasheet for Datasets}\label{appendix:datasheet}

% This appendix documents Claw-SWE-Bench (the full 350-instance set) and its 80-instance Lite companion using the Datasheets for Datasets template of \citet{gebru_datasheets}. Sub-section headings follow the seven-section structure of that template; each question is rendered as a \texttt{\textbackslash paragraph} followed by a short answer drawn from Sections~\ref{sec:clawswebench}--\ref{sec:lite} and Appendices~\ref{appendix:datasheet}--I of this paper. Author-, institution-, and release-specific fields that depend on the camera-ready de-anonymization are marked \textit{[USER FILL]}.

















% Our additional contributions (the merged 350-instance set, the Lite-80 subset selection algorithm, the harness-adapter protocol, evaluation scripts, and figures) are released under MIT.

% \paragraph{Have any third parties imposed IP-based or other restrictions on the data associated with the instances?}
% None known beyond the upstream open-source licenses of the source GitHub repositories.

% \paragraph{Do any export controls or other regulatory restrictions apply to the dataset or to individual instances?}
% No. The dataset consists of public open-source software-engineering issue descriptions and test code.

% \subsection*{B.7 Maintenance}

% \paragraph{Who is supporting/hosting/maintaining the dataset?}
% \textit{[USER FILL: maintainer name(s) and hosting entity; redacted at submission time for double-blind review.]}

% \paragraph{How can the owner/curator/manager of the dataset be contacted (e.g., email address)?}
% \textit{[USER FILL: contact email; redacted at submission time for double-blind review.]}

% \paragraph{Is there an erratum?}
% None at submission time. Any errata identified after release will be published on the dataset's tracking page.

% \paragraph{Will the dataset be updated (e.g., to correct labeling errors, add new instances, delete instances)?}
% We anticipate periodic re-construction of the Lite subset as new models or claws become available for the calibration pool, since the Lite parity evidence is fitted against a fixed set of calibration columns; the full 350-instance set will be held stable except for errata corrections. Major versions of either set will be tagged so that prior published numbers remain reproducible.

% \paragraph{If others want to extend/augment/build on/contribute to the dataset, is there a mechanism for them to do so?}
% Yes. Pull requests and issues against the published GitHub repository will be the primary contribution mechanism, in line with standard open-source benchmark practice. Contributors should preserve the harness-adapter protocol described in Section~\ref{sec:construction} so that any added instances are evaluable under the same fixed run-time parameters as the existing 350.

% \paragraph{If the dataset relates to people, are there applicable limits on the retention of the data associated with the instances?}
% Not applicable. The dataset describes public software-engineering artifacts and does not retain personal data beyond what is already public on the source GitHub repositories.

% \paragraph{Will older versions of the dataset continue to be supported/hosted/maintained?}
% Yes. Tagged versions will remain accessible through the GitHub repository so that prior published evaluation numbers remain reproducible; deprecation will be announced in advance via the tracking page if it ever becomes necessary.
% %